\section{Poets of the Reunion}

Here is Chesterton\footnote{\url{http://www.chesterton.org/discover-chesterton/selected-works/the-critic/george-macdonald/}}:

\begin{quotex}
The originality of George MacDonald has also a historical significance, which perhaps can best be estimated by comparing him with his great countryman Carlyle. It is a measure of the very real power and even popularity of Puritanism in Scotland that Carlyle never lost the Puritan mood even when he lost the whole of the Puritan theology. If an escape from the bias of environment be the test of originality, Carlyle never completely escaped, and George MacDonald did. He evolved out of his own mystical meditations a complete alternative theology leading to a completely contrary mood. And in those mystical meditations he learned secrets far beyond the mere extension of Puritan indignation to ethics and politics. For in the real genius of Carlyle there was a touch of the bully, and wherever there is an element of bullying there is an element of platitude, of reiteration and repeated orders. Carlyle could never have said anything so subtle and simple as MacDonald's saying that God is easy to please and hard to satisfy. Carlyle was too obviously occupied with insisting that God was hard to satisfy; just as some optimists are doubtless too much occupied with insisting that He is easy to please. In other words, MacDonald had made for himself a sort of spiritual environment, a space and transparency of mystical light, which was quite exceptional in his national and denominational environment. He said things that were like the Cavalier mystics, like the Catholic saints, sometimes perhaps like the Platonists or the Swedenborgians, but not in the least like the Calvinists, even as Calvinism remained in a man like Carlyle. And when he comes to be more carefully studied as a mystic, as I think he will be when people discover the possibility of collecting jewels scattered in a rather irregular setting, it will be found, I fancy, that he stands for a rather important turning-point in the history of Christendom, as representing the particular Christian nation of the Scots. As Protestants speak of the morning stars of the Reformation, we may be allowed to note such names here and there as morning stars of the Reunion.
\end{quotex}


Now what interests us here is not the ``alternative'' dimensions of MacDonald's theology which in any case were ``alternative'' relative to Calvinism, but rather the extent to which he rediscovered the Tradition on his own. Take his \emph{Phantastes}, for instance, the volume which yelled at Lewis in the bookstall, until he purchased it and had his imagination baptized.

\emph{Phantastes} has a decidedly ``romantic'' theme, in that a man is baptized Anodos by a bequest from his dead father and his fey grandmother, and sent into Faerie Land in quest of ``he knows not what''. This takes a turn toward the ``eternal feminine'', as various female figures either protect or assault him in his progress on the way, in which he meets a knight with rust-colored armor, who warns him against the a particular lady. Indeed, MacDonald seems to have entirely spiritualized the erotic in this fantasy; we would not follow the modern turn in seeing this as merely an underhanded sublimation, but would rather point towards (say) Mouravieff's doctrine of ``polar beings''.

Among other remarks, MacDonald has his character wonder aloud if ``all men have souls''. At one point in the story, Anodos goes to the ``Church of Darkness'' where the ogress is reading the Scripture, explaining how darkness is not merely coeval, but dominant, over Light. It is there that he finds his shadow, which attaches itself to him, and renders him incapable of rightly discerning Faerie Land.

\begin{quotex}

So, then, as darkness had no beginning, neither will it ever have an end. So, then, it is eternal. The negation of aught else, is its affirmation. Where the light cannot come, there abideth darkness. The light doth but hollow a mine out of the infinite extension of darkness. And ever upon the steps of light treadeth the darkness; yea, springeth in fountains and wells amidst it, from the secret channels of its mighty sea. Truly, man is but a passing flame, moving unquietly amid the surrounding rest of night; without which he yet could not be, and whereof he is in part compounded.

\end{quotex}


Here is one critic who attempts\footnote{\url{http://www.snc.edu/northwind/documents/By\_work/Phantastes/A\_Note\_on\_the\_Structure\_and\_Conclusion\_of\_Phantastes\_-\_John\_Docherty.pdf}} to point out in what way \emph{Phantastes} falls short (not merely artistically but poetically) of Dante's masterpiece. Of this ``falling short'' there is of course no question.

Again, what interests us is how MacDonald was able to work out some of these moral or spiritual truths about the Self on his own. We know that he read Novalis, for instance, and some of the Church Fathers, such as Saint Gregory of Nyssa. In \emph{Phantastes}, he mentions Cornelius Agrippa and Albertus Magnus in a short story within the story about magic, in which (interestingly enough) he suggests that Love is always more powerful than magic, and that all magic of the operative sort meant by the word nowadays implies a sort of danger or risk to the practitioner.

What does he propose for us? MacDonald attempted to make the Good both real and beautiful, which is to say, alluring. However, rather than paint in the lurid colors of esoteric magic, he spoke in the fairy tale. It is in the medium of poetry and fable and ``mythopoeia'' that such deep truths can be safely and innocently communicated, particularly to modern man, who is ``lost'' to direct routes back to God. And this is what makes the Inkling project memorable: that MacDonald began the effort to rebaptize one faculty of man (Imagination) in an attempt to hallow what has so often been a source of evil and seduction, and thereby cause man to ``forget his forgetting'' and to remember the primordial Beauty from which he sprung.

More excerpts from MacDonald's work, and a closer investigation into the mechanics of his project, as well as closer correlation with Tradition, are to follow, but we may take the following as the germ of his philosophy:

\begin{quotex}
All that man sees has to do with man. Worlds cannot be without an intermundane relationship. The community of the centreof all creation suggests an interradiating connection and dependence of the parts.
\end{quotex}

\flrightit{Posted on 2013-05-03 by Logres}

